\documentclass[11pt,a4paper]{article}

% --- Packages ---
\usepackage[margin=.75in]{geometry}
\usepackage[utf8]{inputenc}
\usepackage[T1]{fontenc}
\usepackage{lmodern} % Fix for font not found error
\usepackage{titlesec}
\usepackage{enumitem}
\usepackage{hyperref}
\usepackage{xcolor}
\usepackage{fancyhdr}
\usepackage{comment}
\usepackage{cleveref}
\usepackage[bottom]{footmisc}
\usepackage{graphicx}
\usepackage{longtable}
\usepackage{booktabs}
\usepackage{array}
\usepackage{calc}
\usepackage{etoolbox}

% Standard Pandoc setup
\providecommand{\tightlist}{%
  \setlength{\itemsep}{0pt}\setlength{\parskip}{0pt}}

\crefformat{section}{\S#2#1#3}
\crefformat{subsection}{\S#2#1#3}
\crefformat{subsubsection}{\S#2#1#3}
\crefformat{figure}{#2Figure~#1#3}
\crefformat{footnote}{#2\footnotemark[#1]#3}

\linespread{0.9} % Riduce lo spazio tra le linee

% --- Colors ---
\definecolor{primary}{RGB}{0, 0, 0}
\definecolor{links}{HTML}{0000EE}
\definecolor{blue}{HTML}{26428B}

\hypersetup{
    colorlinks=true,
    linkcolor=links,
    urlcolor=links,
    citecolor=links,
}

% --- Section Formatting ---
\titleformat{\section}
  {\normalfont\Large\bfseries\color{blue}}{\thesection}{1em}{}
\titleformat{\subsection}
{\normalfont\large\bfseries\color{blue}}{\thesubsection}{1em}{}
\titleformat{\subsubsection}
{\normalfont\bfseries\color{blue}}{\thesubsubsection}{1em}{}

% --- Header and Footer ---
\pagestyle{fancy}
\fancyhead[L]{\small Giuseppe Capasso}
\fancyhead[R]{\small learn-shell}
\fancyfoot[C]{\thepage}
\renewcommand{\headrulewidth}{0.4pt}

% --- Custom Title Command ---
\newcommand{\proposaltitle}[1]{
    \begin{center}
        \vspace*{0.1cm}
        {\LARGE \bfseries \color{blue} #1} \\[1.5ex]
        {\large \textbf{Giuseppe Capasso}} \\[1ex]
                \small{
            \vspace{0.1cm}
            \textbf{Materia:} learn-shell\\
        }
            \end{center}
    \vspace{0.3cm}
    \hrule height 0.5pt
    \vspace{0.5cm}
}

\begin{document}

\proposaltitle{Learn Linux Shell}

\section{Learn Linux Shell}\label{learn-linux-shell}

\subsection{Anatomy of a Command}\label{anatomy-of-a-command}

A Linux command typically consists of three parts:

\begin{enumerate}
\def\labelenumi{\arabic{enumi}.}
\tightlist
\item
  \textbf{Command}: The action you want to perform.
\item
  \textbf{Options}: Flags or modifiers to change the behavior of the
  command.
\item
  \textbf{Arguments}: Additional information required by the command to
  execute.
\end{enumerate}

Basic syntax:

\begin{Shaded}
\begin{Highlighting}[]
\BuiltInTok{command} \PreprocessorTok{[}\SpecialStringTok{options}\PreprocessorTok{]} \PreprocessorTok{[}\SpecialStringTok{arguments}\PreprocessorTok{]}
\end{Highlighting}
\end{Shaded}

\subsection{Essential Commands}\label{essential-commands}

\subsubsection{\texorpdfstring{\texttt{ls} - List Directory
Contents}{ls - List Directory Contents}}\label{ls---list-directory-contents}

Lists files and directories in the current directory.

\begin{Shaded}
\begin{Highlighting}[]
\FunctionTok{ls} \PreprocessorTok{[}\SpecialStringTok{options}\PreprocessorTok{]} \PreprocessorTok{[}\SpecialStringTok{directory}\PreprocessorTok{]}
\end{Highlighting}
\end{Shaded}

\textbf{Options}: - \texttt{-l}: Use long format, displaying detailed
information (permissions, owner, size, etc.). - \texttt{-a}: Include
hidden files and directories (those starting with \texttt{.}). -
\texttt{-h}: Show human-readable sizes (e.g., KB, MB, GB).

\subsubsection{\texorpdfstring{\texttt{cd} - Change
Directory}{cd - Change Directory}}\label{cd---change-directory}

Changes the current working directory.

\begin{Shaded}
\begin{Highlighting}[]
\BuiltInTok{cd} \PreprocessorTok{[}\SpecialStringTok{directory}\PreprocessorTok{]}
\end{Highlighting}
\end{Shaded}

\subsubsection{\texorpdfstring{\texttt{mkdir} - Create
Directories}{mkdir - Create Directories}}\label{mkdir---create-directories}

Creates new directories.

\begin{Shaded}
\begin{Highlighting}[]
\FunctionTok{mkdir} \PreprocessorTok{[}\SpecialStringTok{directory\_name}\PreprocessorTok{]}
\end{Highlighting}
\end{Shaded}

\subsubsection{\texorpdfstring{\texttt{touch} - Create Empty
Files}{touch - Create Empty Files}}\label{touch---create-empty-files}

Creates empty files or updates the timestamps of existing files.

\begin{Shaded}
\begin{Highlighting}[]
\FunctionTok{touch} \PreprocessorTok{[}\SpecialStringTok{file\_name}\PreprocessorTok{]}
\end{Highlighting}
\end{Shaded}

\subsubsection{\texorpdfstring{\texttt{rm} - Remove Files or
Directories}{rm - Remove Files or Directories}}\label{rm---remove-files-or-directories}

Removes files or directories.

\begin{Shaded}
\begin{Highlighting}[]
\FunctionTok{rm} \PreprocessorTok{[}\SpecialStringTok{options}\PreprocessorTok{]} \PreprocessorTok{[}\SpecialStringTok{file/directory}\PreprocessorTok{]}
\end{Highlighting}
\end{Shaded}

\textbf{Options}: - \texttt{-r}: Recursively remove directories and
their contents. - \texttt{-f}: Force removal without confirmation (use
with extreme caution).

\subsubsection{\texorpdfstring{\texttt{cat} - Concatenate and Display
Files}{cat - Concatenate and Display Files}}\label{cat---concatenate-and-display-files}

Displays the content of files to standard output.

\begin{Shaded}
\begin{Highlighting}[]
\FunctionTok{cat} \PreprocessorTok{[}\SpecialStringTok{file\_name}\PreprocessorTok{]}
\end{Highlighting}
\end{Shaded}

\subsubsection{\texorpdfstring{\texttt{chmod} - Change File
Permissions}{chmod - Change File Permissions}}\label{chmod---change-file-permissions}

Modifies file or directory permissions (read, write, execute).

\begin{Shaded}
\begin{Highlighting}[]
\FunctionTok{chmod} \PreprocessorTok{[}\SpecialStringTok{options}\PreprocessorTok{]}\NormalTok{ mode file}
\end{Highlighting}
\end{Shaded}

\textbf{Options/Modes}: - \texttt{+}: Adds the specified permissions. -
\texttt{-}: Removes the specified permissions. - \texttt{=}: Sets the
permissions exactly as specified.

\subsubsection{\texorpdfstring{\texttt{curl} - Transfer
Data}{curl - Transfer Data}}\label{curl---transfer-data}

Transfers data to or from a server, supporting various protocols (HTTP,
HTTPS, FTP, etc.).

\begin{Shaded}
\begin{Highlighting}[]
\ExtensionTok{curl} \PreprocessorTok{[}\SpecialStringTok{options}\PreprocessorTok{]} \PreprocessorTok{[}\SpecialStringTok{URL}\PreprocessorTok{]}
\end{Highlighting}
\end{Shaded}

\textbf{Options}: - \texttt{-O}: Save the downloaded file with its
original name. - \texttt{-o\ {[}file{]}}: Save the downloaded file with
a specified name. - \texttt{-L}: Follow redirects.

\end{document}
